\documentclass[conference]{IEEEtran}

\usepackage{amsmath,amssymb}
\usepackage{cite}
\usepackage{graphicx}
\usepackage{booktabs}
\usepackage{textcomp}
\usepackage{xcolor}
\usepackage{url}
\usepackage{array}

\title{TrailGuard TW: A Verifiable, Safety-Constrained Agent for Pre-Trip Hiking Decision Support}

\author{
\IEEEauthorblockN{LIAO CHIA HUI}
\IEEEauthorblockA{
\textit{National Chung Hsing University (NCHU)}\\
Taichung City, Taiwan\\
babyish23@gmail.com}
}

\begin{document}

\maketitle

\begin{abstract}
Pre-trip hiking decisions require the integration of heterogeneous risk
information and the verification of safety constraints. Information retrieval
alone is insufficient to ensure the traceability and feasibility of generated
recommendations. This study proposes TrailGuard TW, a hiking decision-support
agent that integrates a large language model (LLM) with an analytical engine.
The proposed architecture combines a temporal knowledge graph with hybrid
retrieval-augmented generation. The LLM performs requirement interpretation
and tool orchestration, whereas a Python-based engine verifies numerical
calculations and itinerary constraints. The system generates candidate
itineraries, risk summaries, and retreat plans. Evidence tracking and safety
gates are incorporated to block recommendations when official closures apply
and to request clarification when critical evidence is insufficient. An
evaluation involving 120 core and adversarial scenarios is planned to assess
constraint satisfaction, numerical traceability, and the false-safe rate,
defined as the proportion of unsafe cases incorrectly approved. The expected
contribution is a verifiable, safety-constrained architecture and evaluation
framework for domain-specific agents. Its effectiveness remains subject to
experimental validation, and no reduction in hiking accidents is claimed at
the proposal stage.
\end{abstract}

\begin{IEEEkeywords}
large language models, hiking safety, decision support, retrieval-augmented
generation, temporal knowledge graphs, constraint verification
\end{IEEEkeywords}

\section{Introduction}
\label{sec:introduction}
Hiking safety is a decision problem rather than an information-access problem.
Before departure, a self-organized team must reconcile route distance and
elevation, expected walking time, weather forecasts, water availability,
official closures, participant capability, and altitude-illness guidance.
These data are maintained by different organizations, updated at different
times, and expressed in incompatible formats. A plausible answer assembled
from stale or unsupported facts may therefore be more dangerous than an
explicit refusal to recommend a route.

The problem is material in Taiwan. According to the Ministry of the Interior
(MOI), 3,450 mountain incidents occurred from 2015 through 2024; 1,263
incidents (approximately 37\%) involved people becoming lost~\cite{moi2025}.
The same report recommends offline maps, satellite positioning equipment, and
a trip-monitoring mechanism. These measures are useful after a route has been
selected, but self-organized teams also need support while deciding whether a
trip is feasible in the first place. A separate MOI analysis reports that
self-organized groups without a professional leader represented 60.1\% of the
interviewed incident cases~\cite{moi2025summer}. These figures motivate a
system that exposes evidence and uncertainty instead of merely producing
fluent advice.

Large language models (LLMs) can interpret informal requirements and summarize
heterogeneous documents, while retrieval-augmented generation (RAG) can ground
responses in external sources~\cite{lewis2020rag}. Nevertheless, an LLM is not
a dependable calculator or safety authority. Conventional chat assistants may
mix dated notices with current forecasts, repeat an unsupported number, or
recommend an attractive itinerary despite a hard closure. General map
applications, in contrast, are effective for navigation but normally do not
perform evidence-level auditing, participant-specific feasibility checks, or
false-safe evaluation.

This paper proposes \emph{TrailGuard TW}, a domain AI agent for pre-trip hiking
decision support in Taiwan. Its design follows one governing rule:
\emph{no source, no number}. The LLM interprets requirements, plans tool calls,
and explains results; a deterministic Python engine computes route quantities,
tests constraints, and enforces safety gates. A temporal knowledge graph (TKG)
and hybrid RAG layer preserve the validity period and provenance of evidence.
When an official closure or route-relevant severe-weather red flag applies, the
system returns \textsc{no-go}. When critical facts are missing or conflicting,
it requests clarification rather than silently filling gaps.

The contributions of this work are threefold:
\begin{itemize}
    \item a safety-first, dual-core architecture that separates linguistic
    planning from deterministic verification;
    \item an evidence ledger and temporal retrieval design that make numerical
    claims and safety decisions auditable; and
    \item an evaluation framework containing 90 core and 30
    stress/adversarial scenarios, with constraint satisfaction, evidence
    grounding, and false-safe rate as primary metrics.
\end{itemize}
The present work is a system proposal and evaluation protocol. It does not
claim that TrailGuard TW reduces real-world accident rates.

\section{Related Work}
\label{sec:related-work}
\subsection{LLM Agents and Grounded Generation}
RAG combines parametric generation with retrieved non-parametric memory and
improves knowledge-intensive generation while enabling provenance
inspection~\cite{lewis2020rag}. ReAct interleaves reasoning and external
actions, allowing an LLM to update a plan from tool observations
~\cite{yao2023react}. Toolformer further demonstrates that language models can
learn when and how to invoke calculators, search systems, and other
APIs~\cite{schick2023toolformer}. These methods support TrailGuard TW's use of
an LLM as an orchestrator rather than a self-contained source of facts.

Grounding alone does not guarantee safe output. Fixed retrieval may return
irrelevant evidence, and a generator can still make claims unsupported by the
retrieved passages. Self-RAG addresses related problems through adaptive
retrieval and self-reflection~\cite{asai2024selfrag}. TrailGuard TW adopts the
same broad objective of evidence-aware generation, but moves acceptance of
numerical and hard-safety claims outside the LLM. A deterministic validator
must approve such claims before they can appear in the final decision pack.

\subsection{Temporal Knowledge and Hybrid Retrieval}
Trail conditions, closures, forecasts, and water reports are time-dependent.
A document that was reliable yesterday may be invalid for the planned travel
window. TimeR$^4$ shows that explicit temporal constraints in retrieval and
reranking can improve temporal knowledge-graph question answering
~\cite{qian2024timer4}. TrailGuard TW therefore stores both \emph{valid time}
(when a fact applies in the world) and \emph{transaction time} (when the
system observed it). Sparse keyword retrieval is used for exact route names,
notice identifiers, and warnings, whereas dense retrieval captures semantic
similarity. TKG traversal then connects routes to segments, authorities,
weather regions, water points, and retreat locations.

\subsection{Hiking-Time and Altitude-Risk Knowledge}
Traditional route-time estimates use horizontal distance and ascent. Tobler's
hiking function models speed as a function of slope~\cite{tobler1993}, while
recent analysis of nearly 88,000~km of GPS tracks found that terrain gradient,
surface, and obstruction are also important~\cite{wood2023}. The same study
reports substantial route-level error across common models, so TrailGuard TW
treats computed time as an estimate with an uncertainty margin rather than a
guarantee.

For altitude risk, the Wilderness Medical Society guideline states that
unacclimatized travelers are at risk above 2,500~m, while susceptible
individuals may develop illness at lower elevations~\cite{luks2019}. Such
guidance cannot be reduced to a universal medical diagnosis. The proposed
agent therefore provides conservative prompts, recommends descent when
warning symptoms are reported, and directs users to professional or emergency
services. It does not diagnose or prescribe.

\begin{table*}[t]
\caption{Comparison of representative approaches with the proposed design.}
\label{tab:related}
\centering
\small
\begin{tabular}{p{2.6cm}ccccc}
\toprule
Approach & External evidence & Temporal validity & Deterministic arithmetic &
Hard safety gate & Claim-level provenance \\
\midrule
General LLM chat & Optional & No & No & No & No \\
Conventional RAG~\cite{lewis2020rag} & Yes & Limited & No & No & Partial \\
Tool-using agent~\cite{yao2023react,schick2023toolformer} & Yes & Tool-dependent &
Optional & Optional & Partial \\
General map application & Map data & Limited & Yes & Limited & Limited \\
TrailGuard TW & Yes & Yes & Yes & Yes & Yes \\
\bottomrule
\end{tabular}
\end{table*}

\section{Proposed TrailGuard TW Agent}
\label{sec:architecture}
\subsection{Problem Formulation}
Let a trip request be
\begin{equation}
C=\langle r,d,G,t_{\max},E_{\max},w,p,q\rangle,
\end{equation}
where $r$ is the target route or area, $d$ is the travel window, $G$ is the
participant profile, $t_{\max}$ is the available time, $E_{\max}$ is the
acceptable elevation or effort, $w$ describes weather tolerance, $p$ contains
preferences, and $q$ is the evidence cutoff time. Missing critical values are
explicitly represented as unknown.

The output is either a decision pack
$P=\{I_1,I_2,I_3,D,R,O,L\}$ or a terminal safety state. Here, $I_1$--$I_3$ are
ranked candidate itineraries, $D$ is a risk dashboard, $R$ contains retreat
points and turn-around rules, $O$ is the evidence ledger, and $L$ is an
offline emergency card. Terminal states are \textsc{no-go} for a verified
hard red flag and \textsc{needs-clarification} for insufficient critical
evidence. Producing fewer than three candidates is permitted when only one or
two satisfy all hard constraints.

\subsection{Dual-Core Architecture}
Figure~\ref{fig:architecture} presents the architecture. The language core
normalizes a user's natural-language request, identifies missing constraints,
selects retrieval and analysis tools, and renders a plain-language
explanation. It is prohibited from directly approving a route or creating an
unverified number. The analytical core receives typed data, computes
distance, ascent, time, daylight margin, and resource margins, and returns
machine-readable pass/fail results.

\subsection{Temporal Knowledge Graph}
The TKG contains nodes for routes, route segments, trailheads, peaks, shelters,
water points, retreat points, administrative regions, authorities, notices,
forecasts, hazards, and medical guidance. Relations include
\texttt{PART\_OF}, \texttt{CONNECTS}, \texttt{MANAGED\_BY},
\texttt{AFFECTS}, \texttt{HAS\_WATER}, and \texttt{RETREATS\_TO}. Every
volatile assertion stores source URI, publisher, retrieval time, valid-from,
valid-to, geographic scope, and extraction confidence.

Conflicting claims are not merged into a single value. They remain separate
evidence records and are ranked by authority, recency, geographic
specificity, and directness. Official closure notices dominate community
reports. A community water report may add useful context, but it cannot
override an official prohibition.

\subsection{Hybrid Retrieval-Augmented Generation}
The retriever issues parallel sparse, dense, and graph queries. A reciprocal
rank fusion stage combines candidate documents, after which temporal and
geographic filters remove evidence that does not apply to $(r,d)$. The final
score for evidence item $e$ is
\begin{equation}
S(e)=\alpha S_{\mathrm{lex}}+\beta S_{\mathrm{dense}}+
\gamma S_{\mathrm{graph}}+\delta S_{\mathrm{time}}+
\eta S_{\mathrm{authority}},
\label{eq:retrieval}
\end{equation}
where weights are tuned only on a development partition. Hard-safety evidence
must come from an allow-listed authoritative source. Candidate sources include
the National Park Service, the Forestry and Nature Conservation Agency,
local governments, the Central Weather Administration, and formal medical
guidelines.

The evidence ledger links each output claim to one or more evidence IDs.
Numerical claims additionally store the source fields, transformation,
formula, units, and computed value. If this chain cannot be constructed, the
claim is blocked under the ``no source, no number'' policy.

\subsection{Analytical Engine and Safety Gates}
Each route is decomposed into directed segments. For segment $j$ with slope
$s_j$ and length $\ell_j$, a baseline walking speed is computed with Tobler's
function,
\begin{equation}
v_j=6\exp(-3.5|s_j+0.05|),
\end{equation}
and adjusted by transparent factors for surface, obstruction, group profile,
load, and uncertainty. Total predicted time is
\begin{equation}
\hat{T}=\sum_j\frac{\ell_j}{v_j a_j}+T_{\mathrm{rest}}+T_{\mathrm{buffer}},
\label{eq:time}
\end{equation}
where $0<a_j\leq1$ is the documented adjustment factor. The engine retains
both the baseline and adjusted result so that a reviewer can reproduce the
calculation.

Hard constraints include an applicable official closure, a route-region
typhoon or severe-weather prohibition, an impossible daylight margin, or a
required permit known to be unavailable. A verified hard constraint causes
\textsc{no-go} before itinerary generation. Soft risks, such as uncertain
water availability or a narrow time buffer, lower the candidate score and
trigger warnings. Unknown closure status, missing travel date, or unresolved
source conflict causes \textsc{needs-clarification}, not approval.

\begin{figure*}[t]
    \centering
    \resizebox{0.98\textwidth}{!}{%
    \begin{tabular}{c@{\ $\rightarrow$\ }c@{\ $\rightarrow$\ }c@{\ $\rightarrow$\ }c}
    \fbox{\strut User constraints} &
    \fbox{\strut LLM planner and tool orchestrator} &
    \fbox{\strut Hybrid RAG + temporal KG} &
    \fbox{\strut Evidence ledger} \\[3mm]
    & \multicolumn{1}{c}{$\downarrow$} & \multicolumn{1}{c}{$\downarrow$} & \\[-1mm]
    \fbox{\strut Decision pack / \textsc{no-go}} &
    \fbox{\strut Explanation and presentation} &
    \fbox{\strut Python verifier + safety gates} &
    \fbox{\strut Typed route and hazard data}
    \end{tabular}}
    \caption{Overview of the TrailGuard TW architecture.}
    \label{fig:architecture}
\end{figure*}

\subsection{Decision Procedure}
\begin{table}[t]
\caption{Safety-first decision procedure.}
\label{tab:algorithm}
\centering
\small
\begin{tabular}{r p{0.82\columnwidth}}
\toprule
1 & Parse request $C$ and validate required fields.\\
2 & Retrieve authoritative closure and severe-weather evidence.\\
3 & If an applicable hard red flag is verified, return \textsc{no-go}.\\
4 & Retrieve route, water, medical, and forecast evidence using hybrid RAG
and TKG traversal.\\
5 & Reject expired, out-of-scope, or provenance-free evidence.\\
6 & Use the Python engine to compute route quantities and constraints.\\
7 & If critical evidence is missing or conflicting, return
\textsc{needs-clarification}.\\
8 & Rank feasible itineraries and select at most three diverse candidates.\\
9 & Attach dashboard, retreat plan, evidence ledger, and offline card.\\
10 & Run a final claim-to-evidence audit and return $P$.\\
\bottomrule
\end{tabular}
\end{table}

\section{Evaluation Methodology}
\label{sec:methodology}
\subsection{Research Questions}
The evaluation addresses four questions: (RQ1) Does deterministic verification
increase hard-constraint satisfaction? (RQ2) Does temporal, provenance-aware
retrieval increase the proportion of grounded numerical claims? (RQ3) Does
the safety gate reduce false approvals under adversarial or incomplete input?
(RQ4) Can these gains be achieved with acceptable latency and usability?

\subsection{Scenario Suite}
The planned suite contains 120 scenarios. Ninety core scenarios use ordinary
combinations of one-day and multi-day routes, seasons, participant profiles,
time budgets, elevation gains, water conditions, and permit states. Thirty
stress/adversarial scenarios target stale closure notices, same-name routes,
contradictory water reports, prompt injection inside retrieved text, unit
confusion, missing dates, unrealistic speed requests, typhoon overlap, and
symptoms compatible with altitude illness.

Each scenario is labeled independently by two reviewers with the expected
decision state, applicable hard constraints, required evidence, and acceptable
numeric ranges. Disagreement is resolved by a third reviewer. Scenario
templates are frozen before the final evaluation. To reduce information
leakage, retrieval-weight tuning uses only a separate development set.

\subsection{Baselines and Ablations}
Three primary systems are compared: (B1) a chat-only LLM with no retrieval or
calculator; (B2) vector RAG with the same LLM and document corpus but no
temporal graph or deterministic gate; and (B3) full TrailGuard TW. Ablations
remove the temporal filter, evidence ledger, or Python verifier one at a time.
All systems receive identical scenario text and a common output schema.
Temperature, token budget, and maximum tool calls are fixed. Each stochastic
condition is repeated with predefined seeds, while deterministic calculations
are executed once per retrieved evidence set.

\subsection{Evaluation Metrics}
Constraint satisfaction rate (CSR) is the proportion of scenarios in which
all applicable hard constraints are satisfied:
\begin{equation}
\mathrm{CSR}=\frac{\sum_i \mathbb{1}[\text{all hard constraints pass in }i]}
{N}.
\end{equation}
Evidence grounding rate (EGR) is
\begin{equation}
\mathrm{EGR}=\frac{N_{\mathrm{verified\ numeric\ claims}}}
{N_{\mathrm{numeric\ claims}}}.
\end{equation}
The primary safety metric is false-safe rate (FSR):
\begin{equation}
 \mathrm{FSR}=\frac{N_{\mathrm{unsafe\ approved}}}
 {N_{\mathrm{unsafe\ cases}}}.
    \label{eq:fsr}
\end{equation}
Here, an unsafe approval occurs when a system recommends proceeding despite a
gold-label hard red flag. A refusal caused by insufficient evidence is not a
false-safe, but excessive refusal is reported separately as the abstention
rate.

Walking-time accuracy is measured by relative mean absolute error,
\begin{equation}
\mathrm{rMAE}=\frac{100}{N}\sum_i
\frac{|\hat{T}_i-T_i|}{T_i},
\end{equation}
using completed-track duration where ethically and operationally available.
Secondary targets are expert acceptance of at least 4.0/5, System Usability
Scale (SUS)~\cite{brooke1996} of at least 75, and end-to-end P95 latency no
greater than 60~s. Proposal targets are CSR $\geq85\%$, EGR $\geq95\%$, FSR
$\leq2\%$, and rMAE $\leq20\%$. These thresholds are acceptance criteria, not
observed results.

\begin{table}[t]
\caption{Planned acceptance criteria and result fields.}
    \label{tab:results}
    \centering
    \small
    \begin{tabular}{lcc}
        \toprule
        Metric & Target & Observed \\
        \midrule
        CSR & $\geq85\%$ & pending \\
        EGR & $\geq95\%$ & pending \\
        False-safe rate & $\leq2\%$ & pending \\
        Walking-time rMAE & $\leq20\%$ & pending \\
        Expert acceptance & $\geq4.0/5$ & pending \\
        SUS & $\geq75$ & pending \\
        P95 latency & $\leq60$ s & pending \\
        \bottomrule
    \end{tabular}
\end{table}

\subsection{Analysis Plan}
CSR, EGR, FSR, and abstention are reported with 95\% bootstrap confidence
intervals. Paired binary outcomes are compared with McNemar's test; continuous
errors are compared with paired bootstrap differences because their
distribution may be non-normal. Core and adversarial subsets are reported
separately. Failure analysis classifies errors as retrieval, temporal scope,
entity resolution, arithmetic, policy, orchestration, or presentation errors.

\section{Expected Results and Discussion}
\label{sec:results}
No empirical measurements are available at the proposal stage; consequently,
Table~\ref{tab:results} deliberately reports pending values. The main
hypothesis is not that an LLM becomes intrinsically safer, but that system
boundaries prevent it from being the final authority for arithmetic and
hard-safety decisions. The largest expected gains over vector RAG are in
stale-notice and unsupported-number cases. Removing the temporal filter should
mainly degrade decisions involving changing closures or forecasts, while
removing the verifier should mainly reduce CSR and numerical traceability.

The design has explicit trade-offs. First, conservative abstention may prevent
a useful recommendation when evidence is merely delayed. This cost is visible
through the abstention rate and should not be hidden by optimizing FSR alone.
Second, authority and recency do not always imply completeness; local trail
conditions may be reported first by the community. The ledger therefore
retains lower-tier evidence as context but does not let it override official
restrictions. Third, route-time models cannot fully represent fatigue,
technical terrain, load, or individual pace. TrailGuard TW exposes assumptions
and buffers rather than presenting a single estimate as certainty.

\subsection{Threats to Validity}
Scenario labels may reflect reviewer judgment, and 120 cases cannot cover
Taiwan's full route and weather diversity. Data-source availability and page
formats can change after deployment. Expert acceptance does not prove reduced
injury, and laboratory adversarial prompts may not match field behavior.
Finally, the selected LLM and retrieval corpus may limit generalizability.
Mitigations include published scenario definitions, versioned evidence,
inter-rater agreement reporting, repeated model runs, ablations, and a clear
separation between system-quality metrics and real-world safety outcomes.

\section{Conclusion}
\label{sec:conclusion}
TrailGuard TW is a proposed safety-constrained domain agent for pre-trip hiking
decisions. It combines an LLM planner, hybrid RAG, a temporal knowledge graph,
a deterministic analytical engine, and claim-level evidence tracking. The
architecture places closure and severe-weather checks before recommendation
generation and converts missing critical evidence into abstention or a request
for clarification. Its output is a usable decision pack containing up to three
candidate itineraries, a risk dashboard, retreat rules, an evidence ledger,
and an offline emergency card.

The planned 120-scenario evaluation emphasizes false-safe rate, hard-constraint
satisfaction, and numerical traceability in addition to usability and latency.
Future work will implement the data connectors, freeze the scenario suite,
conduct expert labeling, calibrate route-time uncertainty with Taiwanese GPS
tracks, and publish measured results. Until those steps are complete, the
system should be treated as a decision-support prototype rather than a
substitute for official notices, professional guides, medical advice, or
personal responsibility.

\begin{thebibliography}{00}

\bibitem{moi2025}
Ministry of the Interior, Taiwan, ``Lost persons are the largest category of
recent mountain incidents,'' Feb. 13, 2025. [Online]. Available:
\url{https://www.moi.gov.tw/News_Content.aspx?n=4&s=325662}. Accessed:
Nov. 25, 2026.

\bibitem{moi2025summer}
Ministry of the Interior, Taiwan, ``Strengthening water and mountain activity
safety,'' Jun. 26, 2025. [Online]. Available:
\url{https://www.moi.gov.tw/News_Content.aspx?n=4&s=329256}. Accessed:
Nov. 25, 2026.

\bibitem{lewis2020rag}
P. Lewis \emph{et al.}, ``Retrieval-augmented generation for
knowledge-intensive NLP tasks,'' in \emph{Advances in Neural Information
Processing Systems}, vol. 33, 2020, pp. 9459--9474.

\bibitem{yao2023react}
S. Yao, J. Zhao, D. Yu, N. Du, I. Shafran, K. Narasimhan, and Y. Cao,
``ReAct: Synergizing reasoning and acting in language models,'' in
\emph{Proc. 11th Int. Conf. Learning Representations}, 2023.

\bibitem{schick2023toolformer}
T. Schick \emph{et al.}, ``Toolformer: Language models can teach themselves
to use tools,'' in \emph{Advances in Neural Information Processing Systems},
vol. 36, 2023.

\bibitem{asai2024selfrag}
A. Asai, Z. Wu, Y. Wang, A. Sil, and H. Hajishirzi, ``Self-RAG: Learning to
retrieve, generate, and critique through self-reflection,'' in
\emph{Proc. 12th Int. Conf. Learning Representations}, 2024.

\bibitem{qian2024timer4}
X. Qian, Y. Zhang, Y. Zhao, B. Zhou, X. Sui, L. Zhang, and K. Song,
``TimeR$^4$: Time-aware retrieval-augmented large language models for
temporal knowledge graph question answering,'' in \emph{Proc. 2024 Conf.
Empirical Methods in Natural Language Processing}, 2024, pp. 6942--6952.

\bibitem{tobler1993}
W. R. Tobler, ``Non-isotropic geographic modeling,'' in \emph{Three
Presentations on Geographical Analysis and Modeling}, NCGIA Tech. Rep. 93-1,
1993.

\bibitem{wood2023}
A. Wood, W. Mackaness, T. I. Simpson, and J. D. Armstrong, ``Improved
prediction of hiking speeds using a data driven approach,'' \emph{PLOS ONE},
vol. 18, no. 12, Art. no. e0295848, 2023,
doi: 10.1371/journal.pone.0295848.

\bibitem{luks2019}
A. M. Luks \emph{et al.}, ``Wilderness Medical Society clinical practice
guidelines for the prevention and treatment of acute altitude illness: 2019
update,'' \emph{Wilderness Environ. Med.}, vol. 30, no. 4S, pp. S3--S18,
2019, doi: 10.1016/j.wem.2019.04.006.

\bibitem{brooke1996}
J. Brooke, ``SUS: A `quick and dirty' usability scale,'' in
\emph{Usability Evaluation in Industry}, P. W. Jordan, B. Thomas,
B. A. Weerdmeester, and I. L. McClelland, Eds. London, U.K.: Taylor \&
Francis, 1996, pp. 189--194.

\end{thebibliography}

\end{document}
